% Kerntemperatuur: a \subchapter-level (H2) heading grouped under the single
% "Bijlagen" chapter-level heading in main.tex.
\subchapter{Kerntemperatuur}

\lettrine{H}{oe} gaar iets moet zijn draait niet om één druk-op-de-knop-getal.
Bacteriën gaan sneller dood bij een hogere temperatuur en trager bij een
lagere, en dat is precies waarom een biefstuk medium-rare veilig kan zijn
terwijl gehakt altijd door en door moet. Hieronder staat hoe dat werkt, met
richtwaarden om op terug te vallen.

\vspace{1em}
\begingroup
\setlength{\fboxsep}{12pt}
\setlength{\fboxrule}{0.8pt}
\noindent\fcolorbox{terracotta}{paper}{%
  \begin{minipage}{\dimexpr\textwidth-2\fboxsep-2\fboxrule\relax}
  \noindent{\scshape\color{terracotta}\addfontfeatures{LetterSpace=16.0}Kwetsbare groepen eerst lezen}\par
  \vspace{4pt}
  Voor zwangere vrouwen, jonge kinderen, ouderen en mensen met een
  verminderde weerstand liggen de regels strenger: een infectie met
  listeria, salmonella of de toxoplasma-parasiet is voor hen veel
  gevaarlijker. Kook voor deze groepen niet met de langere, lagere
  temperaturen verderop, maar altijd door en door, tot minimaal 70\,°C.
  Geen rosé vlees, geen zachtgekookt ei, geen rauwmelkse kaas en geen rauwe
  vis. Was na rauw vlees of rauwe kip ook meteen je handen, mes en
  snijplank.
  \end{minipage}}\par
\endgroup
\vspace{1em}

\blockrule{Temperatuur en tijd}
Bij een hoge temperatuur zijn bacteriën vrijwel meteen dood. Bij een lagere
temperatuur duurt dat langer, maar het werkt net zo goed als je de tijd
aanhoudt die erbij hoort. Die tijd begint pas te lopen zodra de kern de
temperatuur uit de tabel echt heeft bereikt, niet zodra de pan of de oven
daarop staat.

\blockrule{Richttabel}
Drie manieren om hetzelfde resultaat te halen: heel kort op hoge
temperatuur, een minuut op middelhoge temperatuur, of tien minuten op een
lagere temperatuur. Kies wat bij je bereiding past.

{\scriptsize\renewcommand{\arraystretch}{1.3}
\begin{longtable}{@{}p{4.1cm}@{\hspace{0.6em}}p{1.4cm}@{\hspace{0.6em}}p{1.4cm}@{\hspace{0.6em}}p{1.6cm}@{}}
& \hfill 2 sec & \hfill 1 min & \hfill 10 min \\
\multicolumn{4}{@{}l@{}}{\textcolor{terracotta}{\rule{9cm}{0.4pt}}} \\[0.3em]
\endhead
Rund-, kalfs- \& lamsvlees & \hfill 68\,°C & \hfill 65\,°C & \hfill 60,5\,°C \\
Varkensvlees & \hfill 68\,°C & \hfill 65\,°C & \hfill 60,5\,°C \\
Gevogelte (kip \& kalkoen) & \hfill 73\,°C & \hfill 69\,°C & \hfill 64,5\,°C \\
Gehakt \& burgers & \hfill 68\,°C & \hfill 65\,°C & \hfill 60,5\,°C \\
\end{longtable}}

\vspace{0.2em}
{\footnotesize\itshape\color{inkfaint}Gebaseerd op Amerikaanse
vleeskeuringsrichtlijnen\cite{usda-fsis-appendixa,fda-foodcode}.\par}

\blockrule{Vlees versus gehakt}
Bij een heel stuk vlees, zoals een biefstuk of een varkenshaas, zitten
bacteriën alleen aan de buitenkant en gaan ze al dood zodra je het vlees
dichtschroeit in een hete pan. De kerntemperatuur is dan vooral van belang
bij sous-vide garen, of als het vlees geprikt is (getenderiseerd), waardoor
bacteriën naar binnen kunnen zijn gebracht. Bij gehakt en worst is de hele
massa vermalen, dus wat eerst buitenkant was zit nu overal doorheen. Gehakt
moet daarom altijd volledig gaar zijn, nooit rosé.

\blockrule{Meten in de oven}
\begin{steps}
\item Prik de thermometer horizontaal in het dikste deel van het vlees, uit
de buurt van een bot of een vetrand: die geleiden de warmte anders en geven
een verkeerde meting.
\item Zet de oven laag, tussen 90 en 120\,°C, in plaats van op 200\,°C. Zo
loopt de kerntemperatuur rustig op in plaats van in een paar minuten voorbij
de tabel te schieten.
\item Zodra de thermometer de temperatuur uit de tabel aangeeft, start de
tijd die erbij hoort.
\item Loopt de temperatuur ondertussen te hard door, zet de ovendeur dan op
een kier of de oven helemaal uit. Zolang de kern niet onder de temperatuur
zakt, blijft het veilig.
\item Werk je liever snel op hoge temperatuur, haal het vlees dan een paar
graden voor de doeltemperatuur uit de oven: door nagaren loopt de kern op de
snijplank vanzelf nog wat op.
\end{steps}

\vspace{0.6em}
{\footnotesize\itshape\color{inkfaint}Dit zijn richtwaarden, geen officieel
advies. Twijfel je? Kies de veilige kant: langer of heter.\par}
